\section{Related Work}
\label{sec:related_work}

Steganography in natural language processing (NLP) has evolved from simple rule-based methods—such as zero-width characters or spelling variations—to generative approaches that leverage the probabilistic nature of language models. 

\para{Generative Steganography.} Early work in neural linguistic steganography focused on mapping secret bits to tokens according to their predicted probabilities using techniques like arithmetic coding (AC) or Huffman coding (HC) \cite{ziegler2019neural, yang2021near}. These methods often target high-capacity encoding where a signal is embedded in every generated token. More recent "black-box" approaches like \textsc{LLM-Stega} \cite{lu2024llmstega} balance perceptual and statistical imperceptibility by subtly altering the probability distribution during generation.

\para{LLM-specific Steganography.} The advent of Large Language Models has introduced new steganographic paradigms. \citet{meier2025trojanstego} demonstrated how LLMs can be fine-tuned to encode secrets in output tokens using parity of token IDs, while maintaining fluency. This "Trojan" behavior allows for stealthy exfiltration that is triggered by specific prompt patterns. \citet{westphal2026hide} proposed embedding-space steganography, which uses hyperplane projections in the model's internal activations to partition tokens into buckets for secret bits. They argued that mechanistic interpretability (e.g., linear probes) is a superior detection signal compared to traditional output-level metrics like KL divergence or perplexity.

\para{Sparse and Multi-modal Steganography.} The concept of steganographic sparsity—where hidden messages are not present in every token—has been explored in the context of multi-interaction dialogues. \citet{huang2025gsdfuse} introduced \textsc{GSDFuse}, a framework for detecting "extreme steganographic sparsity" in social media by fusing multi-dimensional weak signals and hierarchical multi-modal features. This approach addresses signals that are dispersed across multiple contexts rather than concentrated in a single sentence.

\para{Positioning.} Unlike previous work that focuses on either high-capacity generative methods or model-internal Trojan behaviors, our work systematically investigates the "dilution limit" for detection and extraction in pure text. We focus on {\bf trigger-based dilution}, a scenario where the presence of a signal is tied to specific, rare linguistic triggers. This allows us to quantify the point where a signal becomes indistinguishable from noise for a general-purpose LLM, even when the model is explicitly prompted to search for it.
